\subsection{R 的子集的上确界}

定理 3. 实直线的每个非空的上有界（相应地下有界）子集都有上确界（相应地下确界）。

设 A 是 R 的非空上有界子集，b 是 A 的一个上界，于是 \(A \subset ] \leftarrow, b]\) 。对每个 \(x \in A\) ，考察集 \(A_x\) ：它由属于 A 且 \(\geqslant x\) 的数组成；诸集 \(A_x\) 在 R 上组成一个滤子基 B，因为 \(A_y \subset A_x\) （若 \(y \geqslant x\)）。设 a 是 A 的一点。对每个属于 A 的 \(x \geqslant a\) ，\(A_x\) 含于紧区间 {[}a, b{]} ，从而滤子基 B 有一个聚点 c。由于诸区间 \([x, \rightarrow[\) 是闭的，c 属于它们的交，因此 c 是 A 的一个上界。但另一方面，A 的每个上界 z 都 \(\geqslant c\) ，否则 c 的邻域 \(]z, \rightarrow[\) 将不含 A 的任何点。于是 c 是 A 的上确界。

对非空的下有界集 B 可类似论证，或者只需注意 -B 非空且上有界，且若 c 是 -B 的上确界，则 -c 是 B 的下确界。

上确界 \(c\) （ \(A\) 的）可由下面两条性质刻画：

\begin{enumerate}
\def\labelenumi{(\roman{enumi})}
\item
  对每个 \(x \in A, x \leqslant c_*\) 。
\item
  对每个 \(a < c\) ，存在 \(x \in \mathbf{A}\) 使得 \(a < x \leqslant c\) 。
\end{enumerate}

闭集（非空且上有界）的上确界属于该集并且是它的最大元素；而 R 的任何非空上有界子集 A 的上确界可以定义为 A 的闭包中的最大实数。

